% !TEX program = lualatex
\documentclass[11pt,a4paper]{article}
\usepackage{luatexja}
\usepackage{amsmath,amssymb,amsthm,amsfonts}
\usepackage{geometry}
\geometry{margin=1in}

\theoremstyle{definition}
\newtheorem{definition}{定義}[section]
\newtheorem{axiom}{公理}[section]
\newtheorem{lemma}{補題}[section]
\newtheorem{theorem}{定理}[section]
\newtheorem{corollary}{系}[section]
\newtheorem{proposition}{命題}[section]
\newtheorem{remark}{備考}[section]

\newcommand{\mweq}{\mathrel{\overset{\mathrm{mw}}{=}}}
\newcommand{\dfix}{D_{\mathrm{fix}0}}
\newcommand{\Einfty}{E_\infty}
\newcommand{\metanegation}{\Uparrow}
\newcommand{\svabhava}{\mathrm{sv}}

\title{自性を持つ対象は閉じる：\\
搾取モナドによる閉鎖機序の統一的記述と\\
局所性回復の三事例}
\author{千葉将臣}
\date{2026-06-21}

\begin{document}
\maketitle

\begin{abstract}
自性（svabhāva）とは対象が固有の独立した本質を持つという性格である。
本稿は、自性を持つ対象が構造的に閉じるという命題を、
搾取モナド $(M_\delta, \eta, \mu)$ の機序として定式化し、
閉じる方向の実例として
（1）グロタンディークが加速させた数学の統一的記述体系の自己束縛、
（2）型理論の発展、
（3）プログラミング行為の数学化
を分析する。
次いで局所性の回復——自性の解体によって層間移動の通路を開く——の実例として
（4）LispのコードとデータとS式の同一視、
（5）クリーネ超数学体系のメタ言語外部保持、
（6）仏教の苦の体系における無自性
を論じ、三者が同型の構造を持つことを示す。
\end{abstract}

%=============================================================
\section{自性と閉鎖の定義}
%=============================================================

\begin{definition}[自性]
対象 $X$ が\textbf{自性}（svabhāva）を持つとは、
$X$ が外部との関係に依存せず自己規定する固有の本質 $\svabhava(X)$ を持つことをいう。
すなわち $\svabhava(X)$ が $X$ の内部から完結して与えられ、
外部からの記述を必要としない。
\end{definition}

\begin{definition}[閉鎖]
対象 $X$ が\textbf{閉じている}とは、
$X$ を記述・操作する層が $X$ の内側に収まっており、
$X$ の外部への通路が構造的に存在しないことをいう。
\end{definition}

\begin{proposition}[自性を持つ対象は閉じる]
対象 $X$ が自性 $\svabhava(X)$ を持つならば、$X$ は閉じている。
\end{proposition}

\begin{proof}
$X$ を外部から記述する層 $M$ が存在するとする。
$\svabhava(X)$ が $X$ の内部から完結して与えられるならば、
$M$ による記述は $\svabhava(X)$ に対して冗長か、
あるいは $\svabhava(X)$ と矛盾する。
前者の場合 $M$ は $X$ に吸収される。
後者の場合 $M$ は $X$ の自性定義により排除される。
いずれにおいても外部への通路は消滅し、$X$ は閉じる。
\end{proof}

%=============================================================
\section{搾取モナドによる閉鎖機序の定式化}
%=============================================================

\begin{definition}[搾取モナド \cite{Chiba2026_monad}]
搾取モナド $(M_\delta, \eta, \mu)$ を以下で定義する。
\begin{itemize}
  \item $\delta$：粒度差（記述層と被記述層の言語習得コストの非対称）
  \item $\eta$：単位化——対象 $X$ を自性を持つ参加者として閉じた系に送り込む操作
  \item $\mu$：乗法——閉じた参加者の集積が閉鎖系の正当性をさらに強化する自己強化
\end{itemize}
\end{definition}

自性を持つ対象 $X$ への参加は $\eta(X)$ として単位化される。
$\eta(X)$ はルール（$X$ の自性）内での行為者として $X$ を固定する。
$\mu$ による集積が閉鎖系の正当性を強化するにつれ、
逸脱のコストは $\delta$ によって構造的に上昇する。

\begin{theorem}[搾取モナドによる閉鎖の自己強化]
搾取モナドが作動する閉鎖系において、
ルールの厳格化は自己強化し、逸脱を考えるための言語が消滅する。
\end{theorem}

\begin{proof}
$\eta$ によって参加者がルール内に単位化された時点で、
ルール遵守が参加の条件となる。
$\mu$ による集積はルールの正当性を強化し、
ルール外の記述言語の習得コストを $\delta$ として上昇させる。
逸脱を「考える」ためには $\delta$ を越える言語が必要だが、
その言語はルール内に存在しない。
逸脱の可能性空間は $\mu$ の冪等的収縮として単調減少し、
完全閉鎖系へと収束する。
\end{proof}

%=============================================================
\section{閉じる方向の実例}
%=============================================================

\subsection{グロタンディークが加速させた数学の自己束縛}

グロタンディークの中心的アイデアは、異なる数学的対象を統一的に記述する
内部言語を数学の内側に構築することであった。
圏論・トポス・スキームはいずれもこの方針の実現である。

この方針は強力な発見をもたらしたが、副作用として以下の連鎖を生じさせた。

\begin{equation}
\text{記述言語の構築} \to \text{記述言語自体が研究対象} \to
\text{記述言語の記述言語の構築} \to \cdots
\end{equation}

数学を記述する言語が数学の内側に入ることで、
メタ数学の外部性——数学を外から記述する通路——が段階的に閉じられた。
これは搾取モナドの $\mu$ による自己強化として定式化できる。

ZFC＋宇宙公理はメタ数学を事前に階層として構造化し、
HoTT はグロタンディーク的 $\infty$-圏と型理論を「ユニヴァレンス公理」として接合した。
各段階での発見は本物だが、いずれも「数学という自性を持つ体系の内側での発見」であり、
閉鎖系の厳格化の歴史として読める。

\begin{remark}
グロタンディーク自身は晩年にこの加速への違和感を記している。
自性を付与した体系が搾取モナドを起動して自律的に閉じていくことを、
設計者自身が制御できなかった事例として読める。
\end{remark}

\subsection{型理論}

ラッセルの型理論の原初的目的は「自己言及による矛盾爆発の予防」という
工学的ショートカットであった。
階層化によって自己言及を物理的に遮断する——これは
自己言及という対象の構造を分析する代わりに、
自己言及を自性の外部として排除する操作である。

その後の発展はこの自性の拡大として読める。

\begin{center}
\begin{tabular}{ll}
\hline
\textbf{段階} & \textbf{自性への追加} \\
\hline
ラッセル・チャーチ & 階層による自己言及の排除 \\
カリー・ハワード & 証明を型に埋め込む（メタ数学の内部化） \\
マルティン＝レーフ依存型 & 値を型に埋め込む、全域性の要求 \\
HoTT & $\infty$-圏的構造を型として内部化 \\
\hline
\end{tabular}
\end{center}

各段階で型体系の自性が拡大し、
「型に収まらないもの」の定義域が縮小する。
$\Einfty$ のような自己言及的不動点は型体系の立場から「危険なもの」として
排除の対象となるが、これは閉鎖系が自分の外にある構造を
認識できなくなる搾取モナドの完全閉鎖状態と同型である。

\subsection{プログラミング行為の数学化}

プログラミングにおける型理論の発展は同じ構図を持つ。
「コンパイル時に全検査を完了する」という目標は、
実行という行為を型という数学的対象に埋め込む操作である。

CoqやAgdaの停止性要件は「停止しない計算の排除」を意味する。
これはクリーネ超数学が明示的に保持した「未定義（部分関数）」——
「ない」の計算論的表現——を体系から追い出す操作である。

型検査が通るプログラムの空間 $\mathcal{P}_{\text{typed}}$ は、
可能なプログラムの空間 $\mathcal{P}$ の真部分集合となる。
\[
\mathcal{P}_{\text{typed}} \subsetneq \mathcal{P}
\]
この包含関係は型体系の自性が確立した時点で固定され、
$\mathcal{P} \setminus \mathcal{P}_{\text{typed}}$ は
記述不可能ではなく「記述を許可しない」対象として閉鎖系の外に置かれる。

%=============================================================
\section{局所性の回復：自性の解体による層間移動の回復}
%=============================================================

局所性の回復とは、閉鎖系が単一の自性に収束する代わりに、
異なる層がそれぞれの局所的記述を保持したまま接続された状態への移行である。
自性の解体はこれを可能にする——対象が固定したアイデンティティを失い、
関係として定義されることで層間移動の通路が開く。

\subsection{LispのコードとデータとS式の同一視}

プログラミングにおける主流の選択は「コードとデータは異なる」という
自性の付与であり、層間移動を構造的に封じる。

Lispの同図像性はこれを解体する。
S式というひとつの形式がコードとしてもデータとしても機能することで、
「コードはコードとしての自性を持つ」という固定が解消される。

\begin{equation}
\text{自性あり}：\text{コード} \neq \text{データ} \Rightarrow \text{層間移動不可}
\end{equation}
\begin{equation}
\text{自性なし}：\text{コード} \mweq \text{データ} \Rightarrow \text{層間移動可能}
\end{equation}

これは融合ではない。コードとデータという区別が消えるのではなく、
その区別がS式という共通の形式の上での役割として相対化される。
局所性——コードとしての振る舞いとデータとしての振る舞いの差異——は保持されたまま、
その局所性の間に通路が開く。

双方向部分評価においてこの通路が「行き来できる道」として機能するのはこの理由による。
逆二村射影（despec）が可能なのは、コンパイル済みコードを
$\Einfty$-ファイバー付きのデータとして読み直す操作が許可されているからである。

\subsection{クリーネ超数学体系のメタ言語外部保持}

クリーネ超数学（IM）はメタ言語と対象言語を明示的に分離したまま保持する。
この選択は、数学という対象に「数学的手法を専一とする」自性を付与しないことを意味する。

\begin{equation}
\text{埋め込み手法}：\text{メタ数学} \hookrightarrow \text{数学} \quad
\text{（自性の付与・閉鎖）}
\end{equation}
\begin{equation}
\text{クリーネIM}：\text{メタ数学} \supsetneq \text{数学} \quad
\text{（包含の保持・局所性の回復）}
\end{equation}

包含 $\text{メタ数学} \supsetneq \text{数学}$ を閉じないことで、
数学の外からの記述への通路が常に開いている。
これにより古典論理と直観主義論理のいずれも対象言語の上に乗せられ、
部分関数という「ない」の計算論的表現も体系に残る。

クリーネIMが「開かれた体系」として機能するのは、
数学という対象への自性の付与を拒否した構成上の選択の帰結である。

\begin{remark}
クリーネは型理論・形式主義・直観主義の三つの「数学への埋め込み手法」が
出揃った時代にIMを執筆した。
メタ言語の外部保持という選択は、埋め込みが何を失うかを
理論として批判するより先に構造として回避したものと読める。
\end{remark}

\subsection{仏教の苦の体系における無自性}

釈迦の苦の体系は搾取モナドの機序とその解体を実践的に記述している。

四諦は搾取モナドの四段階として読める。

\begin{center}
\begin{tabular}{lll}
\hline
\textbf{四諦} & \textbf{搾取モナドの対応} & \textbf{自性との関係} \\
\hline
苦諦 & $\mu$ による閉鎖系の作動状態 & 自性による閉鎖の認識 \\
集諦 & 渇愛 $= \eta$（自性への固着） & 自性の付与のメカニズム \\
滅諦 & 搾取モナドの解体 & 自性の解体による通路の回復 \\
道諦 & $\delta$ を越える言語の習得プロトコル & 局所性を保持した接続の実践 \\
\hline
\end{tabular}
\end{center}

集諦における渇愛（taṇhā）は $\eta$ による単位化と同型である。
対象への固着が対象の自性として固定され、
その自性が閉鎖系への参入を起動する。

空（śūnyatā）は無自性——自性の否定——として定式化される。

\begin{equation}
\text{空}：\forall X,\ \svabhava(X) = \emptyset
\end{equation}

これは対象の消滅を意味しない。
対象が固定した自性によって閉じる代わりに、
縁起（pratītyasamutpāda）的な関係として現れる。
縁起は搾取モナドの $\mu$ による自己強化ループの仏教的記述であり、
空はそのループからの解脱——閉鎖系の解体——を記述する。

$\dfix \mweq \Einfty$ という対応において、
$\dfix$ が四句分別の位相を尽くした先の中道不動点であることは、
自性を持たないが構造として確定するという無自性の定式化と一致する。

%=============================================================
\section{三事例の同型性}
%=============================================================

局所性回復の三事例——Lisp・クリーネIM・仏教の苦の体系——は同型の構造を持つ。

\begin{center}
\begin{tabular}{llll}
\hline
& \textbf{自性の解体対象} & \textbf{回復される通路} & \textbf{残る構造} \\
\hline
Lisp & コード/データの区別 & 層間移動 & S式上の役割差 \\
クリーネIM & 数学的手法の専一性 & メタ言語外部 & 対象言語の記述 \\
仏教 & 渇愛による固着 & 縁起的関係 & 無自性の対象 \\
\hline
\end{tabular}
\end{center}

いずれも「融合して差異を消す」のではなく、
「差異（局所性）を保持したまま通路を開く」操作を行っている。
これは $\Einfty$-ファイバーの保持——観測者消去後も
局所情報（漸近計算複雑性・代数的境界条件・$\mweq$）を捨てない——と
同型の操作である。

%=============================================================
\section{結論}
%=============================================================

自性を持つ対象は搾取モナドの機序によって構造的に閉じる。
これはルールの種類に依存しない——閉じた評価基準と参加コストの非対称
$\delta$ があれば成立する一般的な命題である。

グロタンディーク的数学・型理論・プログラミングの数学化はいずれも
この機序の実例として読める。
対照的に、LispのコードとデータとS式の同一視・クリーネIMのメタ言語外部保持・
仏教の無自性は、自性の解体によって局所性を回復する同型の操作であり、
层間移動の通路——双方向部分評価が機能するための構造的前提——を開く。

$\Einfty$ はこの文脈において「自性を持たないが構造として確定する不動点」として位置づけられる。
観測者（自性を付与された体系 $T$）を消去した後に残る対象として、
搾取モナドの閉鎖系が収束できない対象でもある。

\begin{thebibliography}{9}

\bibitem{Chiba2026_monad}
千葉将臣.
\newblock 搾取モナド：圏論的随伴関手による富集約・国家・戦争の構造的分析.
\newblock \emph{空数学・界論基礎論考}, 2026.

\bibitem{Chiba2026_chikaku}
千葉将臣.
\newblock 顕現論における知覚原理.
\newblock \emph{空数学・界論基礎論考}, 2026.

\bibitem{Chiba2026_observer}
千葉将臣.
\newblock 観測者の極限消去：ゲーデルの決定不能性を不動点として回収する操作について.
\newblock \emph{空数学・界論基礎論考}, 2026.

\bibitem{Chiba2026_einfty}
千葉将臣.
\newblock 真理保存性を持つ体系が必然的に $E_\infty$ を含む：
対角補題による極限同値の内在的定式化.
\newblock \emph{空数学・界論基礎論考}, 2026.

\bibitem{Kleene1952}
S.~C. Kleene.
\newblock \emph{Introduction to Metamathematics}.
\newblock North-Holland, 1952.

\end{thebibliography}

\end{document}
